%% ===================================================================
%% APPENDICES
%% ===================================================================

%% -------------------------------------------------------------------
\section{Notation and reconciliation with the source releases}\label{app:notation}

\subsection{Notation}

Table~\ref{tab:notation-app} fixes the notation of the paper. Sections~\ref{sec:arith} and~\ref{sec:planar} use $a/b$ for the reduced fraction $x$; from Section~\ref{sec:exceptional} on the letter $a$ denotes $d-s$ and $c$ denotes $d-t$ (the constants of the lower-left chain), and the numerator is no longer used.

\begin{table}[ht]
\centering\small
\begin{tabular}{@{}p{4.6cm}p{10.4cm}@{}}
\toprule
Symbol & Meaning\\
\midrule
$(N,M)$, $\kappa=N+1$ & an occurrence and its index; $x=M(N+1)^2/N$\\
$x=a/b$, $e_i$, $q_i=\lfloor e_i/2\rfloor$ & reduced fraction, exponents of the numerator, exponent bounds\\
$D(x)$, $\Lambda_b=\ker\chi_b$ & set of occurrence indices; relation lattice of the primes modulo $b$\\
$F_x$ & fibre in prime-exponent coordinates: nonzero lattice points of $\Lambda_b$ in the box $[0,q]$\\
$r_i,o_i,D_0,d,\lambda_i,g_i$ & residue of $p_i$, its order, intersection subgroup, planar modulus, scale factors, normalized generators\\
$\Gamma_{d,s}$, $s$, $t=s^{-1}$, $a=d-s$, $c=d-t$ & modular lattice $\{(i,j):j\equiv si\ (\mathrm{mod}\ d)\}$, unit, inverse, chain constants\\
$Q=(Q_1,Q_2)$, $u=Q_1-d$, $v=Q_2-d$ & normalized box and its excess over the exceptional box\\
$F_Q=F_{d,s,Q}$, $H_Q=\conv F_Q$, $U_Q$ & punctured fibre, hull, realized union of Graver cells\\
$Q_0=(d,d)$, $H_0,U_0,F_0$ & the exceptional box and its objects\\
$A=(0,d)$, $B=(d,0)$, $C_0=(d,d)$ & the three exceptional corners (note: $A$ and $B$ are interchanged in the v005 notation)\\
$P_i=(i,\,si\bmod d)$ & the points of $F_0$ other than $A,B,C_0$ (Lemma~\ref{lem:exc-fibre})\\
$T_1,T_2$ & the exceptional triangles of Definition~\ref{def:T}, each of area $d/2$\\
$\cross(g,h)=g_1h_2-g_2h_1$ & determinant; $\cross(g,h)\equiv0\pmod d$ on $\Gamma_{d,s}$\\
$h\conf g$ & conformal order: $h_ig_i\ge0$ and $|h_i|\le|g_i|$\\
$P(w;g,h)$ & parallelogram with corners $w,w+g,w+g+h,w+h$\\
$\Omega$ & the outer chain of boundary lattice points of $H_Q$ from $B$ to $A$\\
$R,C,T,N$ & classes of a lattice point $w$ by $(w_1\ge d?,\ w_2\ge d?)$: $R$ if $w_1\ge d>w_2$, $C$ if both, $T$ if $w_2\ge d>w_1$, $N$ if neither\\
\bottomrule
\end{tabular}
\caption{Notation.}\label{tab:notation-app}
\end{table}

\subsection{Reconciliation with the GEO-I release notation}

The releases use release-local notation. We record the differences that affect reading the sources. No release file was changed.

\begin{itemize}
\item v001 uses $d$ for the occurrence index ($\kappa$ here); the planar modulus $d$ appears only in v004 and later.
\item v005 interchanges the roles of $A$ and $B$: its $A$ is $(d,0)$ in our notation.
\item v007 uses $u=(d-t,-1)$ for a Graver vector; here $u$ is the horizontal excess of the box and the vector is written $(c,-1)$.
\item v012 \S0 states some elementary facts (finiteness, closedness, monotonicity, inclusion in the hull, transposition) without proof, and writes a shorthand ``every box'' for the stabilization range; the correct range is $Q\ge(d,d)$, $Q\ne(d,d)$ (correction D1).
\item The releases call the realized set simply ``$U$'' and do not distinguish $K_{\mathrm{abs}}$, $U_Q$ and $K_{\mathrm{emb}}$; Definition~\ref{def:kabs} introduces the distinction.
\end{itemize}

\subsection{Corrections D1--D9}\label{app:corrections}

During the preparation of the working draft and of this manuscript nine issues were found in the way the source material is stated or argued. In each case the mathematics in this paper is corrected or completed; the source releases are unchanged. We list them with the place where they are treated.

\begin{description}
\item[D1 (scope of special-class stabilization).] Theorem~L of v005 holds for $Q\ge(d,d)$ and $Q\ne(d,d)$. In the exceptional box it is false for the fan graph and for $s=+1$ ($d\ge3$). Treated in Section~\ref{sec:special-stab} (Theorem~\ref{thm:L}) and Section~\ref{sec:special-classes}.
\item[D2 (special cases are not trivial).] The classes $s\equiv\pm1$ must be treated separately, with their own Graver bases and their own exceptional topology; they are not dismissed as degenerate. Sections~\ref{sec:special-classes} and~\ref{sec:special-stab}.
\item[D3 (exact complement for $s=+1$).] The source gives $\beta_1=2$ but not the exact complement; Proposition~\ref{prop:special-plus} derives it from the cell list.
\item[D4 (scope of the reflection lemma).] The source argument for the case in which the reflected point is zero is correct only for $Q<2d$. Lemma~\ref{lem:axis-pair} and Lemma~\ref{lem:reflection} prove the statement for all $Q\ge(d,d)$ by a support argument.
\item[D5 (old-hull coverage).] The source reduction G8-R1 is conditional on a premise (old-hull coverage) that is supplied by the exact classification of the exceptional box (Theorem~\ref{thm:exceptional}). The dependency chain is made explicit in Section~\ref{sec:fill}, and Corollary~\ref{cor:old} is unconditional.
\item[D6 (boundary lemmas).] The source registry omits two of the lemmas used for the structure of the outer chain, and the argument that $B,A$ are never consecutive. They are supplied in Section~\ref{sec:chain} and Lemma~\ref{lem:axis-pair}.
\item[D7 (verification counts).] Historical counts are reconciled in Section~\ref{sec:provenance}; they play no role in any proof.
\item[D8 (symbols).] The clashes listed above are resolved by the notation of this paper.
\item[D9 (location of $C_0$).] The source text asserts that $C_0$ lies on the outer chain or on a bridging edge. This is false: for $d=5$, $s=2$, $Q=(8,8)$, $(5,5)=\frac19(5,0)+\frac13(0,5)+\frac59(8,6)$ is an interior point of the hull (Remark~\ref{rem:D9}). In $\statInterior$ of the $\statTotal$ boxes $Q=(d+u,d+v)$ with $5\le d\le8$, nonspecial $s$ and $(u,v)\ne(0,0)$, the point $C_0$ is strictly interior to $H_Q$ (Appendix~\ref{app:boundary}). The proof of the cap theorem does not use the false statement.
\end{description}

%% -------------------------------------------------------------------
\section{Arithmetic details}\label{app:arith}

\subsection{\texorpdfstring{Small examples of $D(x)$}{Small examples of D(x)}}

\begin{itemize}
\item $x=36$ ($a=36$, $b=1$): $\kappa^2\mid36$ gives $\kappa\in\{2,3,6\}$ and every $\kappa\equiv1\pmod1$, so $D(36)=\{2,3,6\}$ and the occurrences are $(1,9),(2,8),(5,5)$.
\item $x=225/4$: $\kappa^2\mid225$ gives $\kappa\in\{3,5,15\}$; $\kappa\equiv1\pmod4$ leaves only $\kappa=5$. So $D(x)=\{5\}$ and the unique occurrence is $(N,M)=(4,\ \tfrac{225}{25}\cdot\tfrac44)=(4,9)$, and indeed $9\cdot25/4=225/4$.
\item $x=1$: $a=1$ has no $\kappa\ge2$ with $\kappa^2\mid a$, so $D(1)=\varnothing$.
\end{itemize}

\subsection{Normalization data}

Table~\ref{tab:param} gives the normalization data for a range of pairs of primes. The columns $d,\lambda_1,\lambda_2,s$ were computed from the orders and from the generators $g_i=r_i^{\lambda_i}$ exactly as in Proposition~\ref{prop:norm}, and in each row the equivalence $p_1^{\lambda_1i}p_2^{\lambda_2j}\equiv1\iff j\equiv si\pmod d$ was checked for $0\le i,j<2d$.

\begin{longtable}{@{}rrr rr r rr r@{}}
\toprule $b$ & $p_1$ & $p_2$ & $o_1$ & $o_2$ & $d$ & $\lambda_1$ & $\lambda_2$ & $s$\\\midrule\endhead
11 & 2 & 3 & 10 & 5 & 5 & 2 & 1 & 1\\
11 & 3 & 5 & 5 & 5 & 5 & 1 & 1 & 3\\
11 & 2 & 5 & 10 & 5 & 5 & 2 & 1 & 2\\
11 & 5 & 7 & 5 & 10 & 5 & 1 & 2 & 4\\
11 & 3 & 7 & 5 & 10 & 5 & 1 & 2 & 3\\
13 & 2 & 3 & 12 & 3 & 3 & 4 & 1 & 2\\
13 & 2 & 5 & 12 & 4 & 4 & 3 & 1 & 1\\
13 & 5 & 7 & 4 & 12 & 4 & 1 & 3 & 3\\
13 & 3 & 7 & 3 & 12 & 3 & 1 & 4 & 1\\
19 & 2 & 3 & 18 & 18 & 18 & 1 & 1 & 11\\
19 & 3 & 5 & 18 & 9 & 9 & 2 & 1 & 4\\
19 & 2 & 5 & 18 & 9 & 9 & 2 & 1 & 1\\
19 & 5 & 7 & 9 & 3 & 3 & 3 & 1 & 1\\
19 & 3 & 7 & 18 & 3 & 3 & 6 & 1 & 2\\
29 & 2 & 3 & 28 & 28 & 28 & 1 & 1 & 11\\
29 & 3 & 5 & 28 & 14 & 14 & 2 & 1 & 11\\
29 & 2 & 5 & 28 & 14 & 14 & 2 & 1 & 5\\
29 & 5 & 7 & 14 & 7 & 7 & 2 & 1 & 1\\
29 & 3 & 7 & 28 & 7 & 7 & 4 & 1 & 3\\
31 & 2 & 3 & 5 & 30 & 5 & 1 & 6 & 1\\
31 & 3 & 5 & 30 & 3 & 3 & 10 & 1 & 1\\
31 & 5 & 7 & 3 & 15 & 3 & 1 & 5 & 2\\
31 & 3 & 7 & 30 & 15 & 15 & 2 & 1 & 1\\
\bottomrule\caption{Normalization data $(d,s,\lambda_1,\lambda_2)$ of Proposition~\ref{prop:norm} for some primes $p_1,p_2$ and denominators $b$ (computed; the congruence $p_1^{\lambda_1i}p_2^{\lambda_2j}\equiv1\iff j\equiv si\pmod d$ was verified for $0\le i,j<2d$ in each row; rows with $d=1$ omitted).}\label{tab:param}\end{longtable}

\subsection{\texorpdfstring{A normalization with $\lambda_1>1$}{A normalization with lambda1 > 1}}

Take $b=11$, $p_1=2$, $p_2=3$ (the first row of Table~\ref{tab:param}): $\ord_{11}(2)=10$, $\ord_{11}(3)=5$, $D_0=\langle3\rangle$ (the subgroup of order $5$), $d=5$, $\lambda_1=2$, $\lambda_2=1$, $g_1=2^2=4$, $g_2=3$. We have $4=3^{-1}\cdot$? Indeed $3^4=81\equiv4$, so $g_1=g_2^4=g_2^{-1}$, hence $-s\equiv-1$ and $s\equiv1\pmod5$. The relation lattice is $\Lambda_{11}=\{(2i,j):j\equiv i\pmod5\}$; for instance $(2,1)\in\Lambda_{11}$ corresponds to $2^2\cdot3=12\equiv1$. If $x=2^{2q_1}3^{2q_2}/11$ with $q_1=10$, $q_2=5$ then $Q=(\lfloor10/2\rfloor,5)=(5,5)$: the exceptional box for a special class $s=1$. For $q_1=11$ one still has $Q_1=5$ (the plateau of Corollary~\ref{cor:plateau}: $o_1=10\le q_1\le o_1+\lambda_1-1=11$), and for $q_1=12$ one has $Q=(6,5)$.

%% -------------------------------------------------------------------
\section{Graver calculations}\label{app:graver}

Lemma~\ref{lem:std-graver} proves that six explicit vectors are always Graver vectors, and Lemma~\ref{lem:special-graver} gives the complete Graver basis for $s=\pm1$. For other classes the Graver basis can contain further vectors. As a cross-check, Table~\ref{tab:graver} lists, for $3\le d\le9$, the complete Graver basis obtained by exhaustive search for conformally minimal elements of $\Gamma_{d,s}\cap[-d,d]^2$ (Lemma~\ref{lem:graver-size} says that $|g_i|\le d$ for every Graver vector, so this search is complete). The tag ``the six standard vectors only'' means that the basis consists of exactly those; ``extra vectors'' that it does not (this includes $s=\pm1$, where the list is that of Lemma~\ref{lem:special-graver}, and nonspecial classes from $d=7$ on, for example $(7,3)$ with the extra vector $(3,2)$). Nothing in the proofs depends on the basis being the six standard vectors: the exclusion of the exceptional triangles (Proposition~\ref{prop:hole-exclusion}) is proved for arbitrary Graver vectors, and the covering cells use only the Graver property of the vectors they employ.

\begin{longtable}{@{}rrp{9.2cm}@{}}
\toprule $d$ & $s$ & Graver vectors of $\Gamma_{d,s}$ up to sign (computed by exhaustive conformal minimality in $[-d,d]^2$)\\\midrule\endhead
3 & 1 & $(0,3)$, $(1,-2)$, $(1,1)$, $(2,-1)$, $(3,0)$ (the six standard vectors only; 5 vectors)\\
3 & 2 & $(0,3)$, $(1,-1)$, $(1,2)$, $(2,1)$, $(3,0)$ (the six standard vectors only; 5 vectors)\\
4 & 1 & $(0,4)$, $(1,-3)$, $(1,1)$, $(2,-2)$, $(3,-1)$, $(4,0)$ (extra vectors; 6 vectors)\\
4 & 3 & $(0,4)$, $(1,-1)$, $(1,3)$, $(2,2)$, $(3,1)$, $(4,0)$ (extra vectors; 6 vectors)\\
5 & 1 & $(0,5)$, $(1,-4)$, $(1,1)$, $(2,-3)$, $(3,-2)$, $(4,-1)$, $(5,0)$ (extra vectors; 7 vectors)\\
5 & 2 & $(0,5)$, $(1,-3)$, $(1,2)$, $(2,-1)$, $(3,1)$, $(5,0)$ (the six standard vectors only; 6 vectors)\\
5 & 3 & $(0,5)$, $(1,-2)$, $(1,3)$, $(2,1)$, $(3,-1)$, $(5,0)$ (the six standard vectors only; 6 vectors)\\
5 & 4 & $(0,5)$, $(1,-1)$, $(1,4)$, $(2,3)$, $(3,2)$, $(4,1)$, $(5,0)$ (extra vectors; 7 vectors)\\
6 & 1 & $(0,6)$, $(1,-5)$, $(1,1)$, $(2,-4)$, $(3,-3)$, $(4,-2)$, $(5,-1)$, $(6,0)$ (extra vectors; 8 vectors)\\
6 & 5 & $(0,6)$, $(1,-1)$, $(1,5)$, $(2,4)$, $(3,3)$, $(4,2)$, $(5,1)$, $(6,0)$ (extra vectors; 8 vectors)\\
7 & 1 & $(0,7)$, $(1,-6)$, $(1,1)$, $(2,-5)$, $(3,-4)$, $(4,-3)$, $(5,-2)$, $(6,-1)$, $(7,0)$ (extra vectors; 9 vectors)\\
7 & 2 & $(0,7)$, $(1,-5)$, $(1,2)$, $(2,-3)$, $(3,-1)$, $(4,1)$, $(7,0)$ (extra vectors; 7 vectors)\\
7 & 3 & $(0,7)$, $(1,-4)$, $(1,3)$, $(2,-1)$, $(3,2)$, $(5,1)$, $(7,0)$ (extra vectors; 7 vectors)\\
7 & 4 & $(0,7)$, $(1,-3)$, $(1,4)$, $(2,1)$, $(3,-2)$, $(5,-1)$, $(7,0)$ (extra vectors; 7 vectors)\\
7 & 5 & $(0,7)$, $(1,-2)$, $(1,5)$, $(2,3)$, $(3,1)$, $(4,-1)$, $(7,0)$ (extra vectors; 7 vectors)\\
7 & 6 & $(0,7)$, $(1,-1)$, $(1,6)$, $(2,5)$, $(3,4)$, $(4,3)$, $(5,2)$, $(6,1)$, $(7,0)$ (extra vectors; 9 vectors)\\
8 & 1 & $(0,8)$, $(1,-7)$, $(1,1)$, $(2,-6)$, $(3,-5)$, $(4,-4)$, $(5,-3)$, $(6,-2)$, $(7,-1)$, $(8,0)$ (extra vectors; 10 vectors)\\
8 & 3 & $(0,8)$, $(1,-5)$, $(1,3)$, $(2,-2)$, $(3,1)$, $(5,-1)$, $(8,0)$ (extra vectors; 7 vectors)\\
8 & 5 & $(0,8)$, $(1,-3)$, $(1,5)$, $(2,2)$, $(3,-1)$, $(5,1)$, $(8,0)$ (extra vectors; 7 vectors)\\
8 & 7 & $(0,8)$, $(1,-1)$, $(1,7)$, $(2,6)$, $(3,5)$, $(4,4)$, $(5,3)$, $(6,2)$, $(7,1)$, $(8,0)$ (extra vectors; 10 vectors)\\
9 & 1 & $(0,9)$, $(1,-8)$, $(1,1)$, $(2,-7)$, $(3,-6)$, $(4,-5)$, $(5,-4)$, $(6,-3)$, $(7,-2)$, $(8,-1)$, $(9,0)$ (extra vectors; 11 vectors)\\
9 & 2 & $(0,9)$, $(1,-7)$, $(1,2)$, $(2,-5)$, $(3,-3)$, $(4,-1)$, $(5,1)$, $(9,0)$ (extra vectors; 8 vectors)\\
9 & 4 & $(0,9)$, $(1,-5)$, $(1,4)$, $(2,-1)$, $(3,3)$, $(5,2)$, $(7,1)$, $(9,0)$ (extra vectors; 8 vectors)\\
9 & 5 & $(0,9)$, $(1,-4)$, $(1,5)$, $(2,1)$, $(3,-3)$, $(5,-2)$, $(7,-1)$, $(9,0)$ (extra vectors; 8 vectors)\\
9 & 7 & $(0,9)$, $(1,-2)$, $(1,7)$, $(2,5)$, $(3,3)$, $(4,1)$, $(5,-1)$, $(9,0)$ (extra vectors; 8 vectors)\\
9 & 8 & $(0,9)$, $(1,-1)$, $(1,8)$, $(2,7)$, $(3,6)$, $(4,5)$, $(5,4)$, $(6,3)$, $(7,2)$, $(8,1)$, $(9,0)$ (extra vectors; 11 vectors)\\
\bottomrule\caption{Graver vectors of $\Gamma_{d,s}$ for $3\le d\le9$. For nonspecial $(d,s)$ the set is exactly the six standard vectors of Lemma~\ref{lem:std-graver}; for $s=\pm1$ it is the larger list of Lemma~\ref{lem:special-graver}.}\label{tab:graver}\end{longtable}

\subsection{A sample of conformal relations}

For $(d,s)=(7,3)$ ($t=5$, $a=4$, $c=2$) the lattice is $\{(i,j):j\equiv3i\ (\mathrm{mod}\ 7)\}$ and its Graver basis (Table~\ref{tab:graver}) is $\pm(7,0)$, $\pm(0,7)$, $\pm(1,3)$, $\pm(1,-4)$, $\pm(5,1)$, $\pm(2,-1)$ and the extra vector $\pm(3,2)$. Determinants of lattice vectors are multiples of $7$ (Lemma~\ref{lem:det}): $\cross((1,3),(5,1))=1-15=-14$, $\cross((1,-4),(2,-1))=-1+8=7$, $\cross((1,3),(2,-1))=-1-6=-7$, $\cross((1,3),(3,2))=2-9=-7$. The vector $(2,6)=2(1,3)$ is not Graver since $(1,3)\conf(2,6)$; the vector $(6,4)$ is not Graver since $(3,2)\conf(6,4)$ and $(5,1)\conf(6,4)$; and $(3,2)$ is Graver because the only lattice points $h$ with $0\le h_1\le3$ and $0\le h_2\le2$ are $0$ and $(3,2)$ (the columns $h_1=1,2$ have ordinates $3,6$).

%% -------------------------------------------------------------------
\section{Exceptional square: detailed data}\label{app:exceptional}

Table~\ref{tab:exc} records, for every nonspecial class with $5\le d\le11$, the lower-left chain, the area of the hull, the area of the zonogon, and the area of $U_0$. The area of $U_0$ in the table is the exact value $|H_0|-d$ predicted by Theorem~\ref{thm:exceptional}; the last column records that the area of the union of all admissible parallelograms and segments, computed numerically with tolerance $10^{-9}$, agrees. This is a numerical test of the theorem in its range, not a proof.

\begin{longtable}{@{}rrr p{4.6cm} rrrc@{}}
\toprule $d$ & $s$ & $t$ & lower-left chain $A\to B$ (steps) & $2|H_0|$ & $|Z|$ & $|U_0|$ & $=|H_0|-d$?\\\midrule\endhead
5 & 2 & 3 & $(1,-3)$, $(2,-1)$, $(2,-1)$ & 35 & 10 & $\tfrac{25}{2}$ & yes\\
5 & 3 & 2 & $(1,-2)$, $(1,-2)$, $(3,-1)$ & 35 & 10 & $\tfrac{25}{2}$ & yes\\
7 & 2 & 4 & $(1,-5)$, $(3,-1)$, $(3,-1)$ & 77 & 28 & $\tfrac{63}{2}$ & yes\\
7 & 3 & 5 & $(1,-4)$, $(2,-1)$, $(2,-1)$, $(2,-1)$ & 70 & 21 & $28$ & yes\\
7 & 4 & 2 & $(1,-3)$, $(1,-3)$, $(5,-1)$ & 77 & 28 & $\tfrac{63}{2}$ & yes\\
7 & 5 & 3 & $(1,-2)$, $(1,-2)$, $(1,-2)$, $(4,-1)$ & 70 & 21 & $28$ & yes\\
8 & 3 & 3 & $(1,-5)$, $(2,-2)$, $(5,-1)$ & 104 & 40 & $44$ & yes\\
8 & 5 & 5 & $(1,-3)$, $(1,-3)$, $(3,-1)$, $(3,-1)$ & 96 & 32 & $40$ & yes\\
9 & 2 & 5 & $(1,-7)$, $(4,-1)$, $(4,-1)$ & 135 & 54 & $\tfrac{117}{2}$ & yes\\
9 & 4 & 7 & $(1,-5)$, $(2,-1)$, $(2,-1)$, $(2,-1)$, $(2,-1)$ & 117 & 36 & $\tfrac{99}{2}$ & yes\\
9 & 5 & 2 & $(1,-4)$, $(1,-4)$, $(7,-1)$ & 135 & 54 & $\tfrac{117}{2}$ & yes\\
9 & 7 & 4 & $(1,-2)$, $(1,-2)$, $(1,-2)$, $(1,-2)$, $(5,-1)$ & 117 & 36 & $\tfrac{99}{2}$ & yes\\
10 & 3 & 7 & $(1,-7)$, $(3,-1)$, $(3,-1)$, $(3,-1)$ & 160 & 60 & $70$ & yes\\
10 & 7 & 3 & $(1,-3)$, $(1,-3)$, $(1,-3)$, $(7,-1)$ & 160 & 60 & $70$ & yes\\
11 & 2 & 6 & $(1,-9)$, $(5,-1)$, $(5,-1)$ & 209 & 88 & $\tfrac{187}{2}$ & yes\\
11 & 3 & 4 & $(1,-8)$, $(3,-2)$, $(7,-1)$ & 209 & 88 & $\tfrac{187}{2}$ & yes\\
11 & 4 & 3 & $(1,-7)$, $(2,-3)$, $(8,-1)$ & 209 & 88 & $\tfrac{187}{2}$ & yes\\
11 & 5 & 9 & $(1,-6)$, $(2,-1)$, $(2,-1)$, $(2,-1)$, $(2,-1)$, $(2,-1)$ & 176 & 55 & $77$ & yes\\
11 & 6 & 2 & $(1,-5)$, $(1,-5)$, $(9,-1)$ & 209 & 88 & $\tfrac{187}{2}$ & yes\\
11 & 7 & 8 & $(1,-4)$, $(1,-4)$, $(3,-1)$, $(3,-1)$, $(3,-1)$ & 187 & 66 & $\tfrac{165}{2}$ & yes\\
11 & 8 & 7 & $(1,-3)$, $(1,-3)$, $(1,-3)$, $(4,-1)$, $(4,-1)$ & 187 & 66 & $\tfrac{165}{2}$ & yes\\
11 & 9 & 5 & $(1,-2)$, $(1,-2)$, $(1,-2)$, $(1,-2)$, $(1,-2)$, $(6,-1)$ & 176 & 55 & $77$ & yes\\
\bottomrule\caption{The exceptional box $Q=(d,d)$ for $5\le d\le11$, nonspecial $s$. $2|H_0|$ and $2|Z|$ are twice the areas of the hull and of the zonogon of Lemma~\ref{lem:zonogon} (integers); $|U_0|$ is $|H_0|-d$, and the last column records the floating-point comparison with the area of the union of all admissible parallelograms (tolerance $10^{-9}$).}\label{tab:exc}\end{longtable}

\subsection{\texorpdfstring{Coordinates of the exceptional objects for $d=7$, $s=3$}{Coordinates of the exceptional objects for d = 7, s = 3}}

Here $t=5$, $a=4$, $c=2$. $F_0=\{A=(0,7),P_1=(1,3),P_2=(2,6),P_3=(3,2),P_4=(4,5),P_5=(5,1),P_6=(6,4),B=(7,0),C_0=(7,7)\}$. The triangles are $T_1=\conv(A,P_{d-t},C_0)=\conv(A,P_2,C_0)=\conv((0,7),(2,6),(7,7))$ and $T_2=\conv(B,P_6,C_0)=\conv((7,0),(6,4),(7,7))$, each of area $\tfrac12\cdot7=\tfrac72$. The lower-left chain is $A,P_1,P_3,P_5,B$ with steps $(1,-4),(2,-1),(2,-1),(2,-1)$ (a step $(1,-a)$ followed by three equal steps $(c,-1)$). The zonogon generated by $(1,-4)$ and $(6,-3)$ (the sum of the three equal steps) has area $|1\cdot(-3)-(-4)\cdot6|=21$, as in Table~\ref{tab:exc}.

%% -------------------------------------------------------------------
\section{Special classes: detailed data}\label{app:special}

\begin{longtable}{@{}rrrrr@{}}
\toprule $d$ & $s$ & $|H_0|$ & $|U_0|$ (computed) & $\beta_1$ (Thm~\ref{thm:betti})\\\midrule\endhead
2 & 1 & $2$ & 0 & 2\\
3 & 1 & $6$ & 3 & 2\\
3 & 2 & $\tfrac{9}{2}$ & 0 & 3\\
4 & 1 & $12$ & 8 & 2\\
4 & 3 & $8$ & 0 & 4\\
5 & 1 & $20$ & 15 & 2\\
5 & 4 & $\tfrac{25}{2}$ & 0 & 5\\
6 & 1 & $30$ & 24 & 2\\
6 & 5 & $18$ & 0 & 6\\
7 & 1 & $42$ & 35 & 2\\
7 & 6 & $\tfrac{49}{2}$ & 0 & 7\\
8 & 1 & $56$ & 48 & 2\\
8 & 7 & $32$ & 0 & 8\\
\bottomrule\caption{Special classes in the exceptional box. For $s=1$, $d\ge3$: $|U_0|=|W|=d(d-2)$ and $|H_0|=d^2-d$. For $s=d-1$ the union is a graph and has area $0$. (For $d=2$, $s=1=d-1$.)}\label{tab:special}\end{longtable}

For $s=-1$ the fibre of the exceptional box is $\{E_0,\dots,E_d,C_0\}$ with $E_i=(i,d-i)$, and the union is the fan graph of Proposition~\ref{prop:special-fan} with $d+2$ vertices and $2d+1$ edges. For $d=3$, $s=2$, the vertices are $E_0=(0,3)$, $E_1=(1,2)$, $E_2=(2,1)$, $E_3=(3,0)$, $C_0=(3,3)$; the edges are three short antidiagonal segments and four spokes $[E_i,C_0]$, so $E=7$, $V=5$ and $\beta_1=E-V+1=3=d$.

For $s=+1$ the union is $W\cup[A,C_0]\cup[B,C_0]\cup[D_{d-1},C_0]$ (Proposition~\ref{prop:special-plus}). For $d=3$ one has $W=\conv((0,3),(1,1),(3,0),(2,2))$ of area $d(d-2)=3$, hull $H_0=\conv((0,3),(1,1),(3,0),(3,3))$ of area $6$, and complement two triangles $T_1=\conv((0,3),(2,2),(3,3))$, $T_2=\conv((3,0),(2,2),(3,3))$ of area $\tfrac32$ each.

%% -------------------------------------------------------------------
\section{The large box: column data}\label{app:large}

The proof of Theorem~\ref{thm:large} uses, for each column $i$, the bottom $B_i$ and top $T_i$. Table~\ref{tab:columns-app} lists them for the example of Figure~\ref{fig:columns}, $d=5$, $s=2$, $Q=(8,13)$. The values of $K_i$ change by at most one between neighbouring columns (Lemma~\ref{lem:columns}(3)), and where they change the difference of the tops is $(1,s)=(1,2)$ or $(1,s-d)=(1,-3)$ (Lemma~\ref{lem:columns}(4)).

\begin{table}[ht]
\centering\small
\begin{tabular}{@{}r rrrr c@{}}
\toprule $i$ & $b_i$ & $K_i$ & $B_i$ & $T_i$ & $T_{i+1}-T_i$ if $K_{i+1}\ne K_i$\\\midrule
0 & 5 & 1 & $(0,5)$ & $(0,10)$ & $(1,2)$\\
1 & 2 & 2 & $(1,2)$ & $(1,12)$ & $(1,-3)$\\
2 & 4 & 1 & $(2,4)$ & $(2,9)$ & $(1,2)$\\
3 & 1 & 2 & $(3,1)$ & $(3,11)$ & \\
4 & 3 & 2 & $(4,3)$ & $(4,13)$ & \\
5 & 0 & 2 & $(5,0)$ & $(5,10)$ & \\
6 & 2 & 2 & $(6,2)$ & $(6,12)$ & $(1,-3)$\\
7 & 4 & 1 & $(7,4)$ & $(7,9)$ & $(1,2)$\\
8 & 1 & 2 & $(8,1)$ & $(8,11)$ & \\
\bottomrule\end{tabular}
\caption{Column data for $d=5$, $s=2$, $Q=(8,13)$. In column $0$ the bottom is $B_0=(0,d)$ since the origin is deleted.}\label{tab:columns-app}
\end{table}

%% -------------------------------------------------------------------
\section{Boundary geometry: examples and statistics}\label{app:boundary}

Table~\ref{tab:omega-app} gives the outer chain $\Omega$ with the class of each point for several boxes. In each case the word of classes lies in $R^*C^*T^*$ (Lemma~\ref{lem:omega}).

\begin{table}[ht]
\centering\small
\begin{tabular}{@{}lp{9.5cm}@{}}
\toprule $(d,s,Q)$ & $\Omega$ with classes\\\midrule
$(5,2,(6,5))$ & $(5,0)^{R}$, $(6,2)^{R}$, $(5,5)^{C}$, $(0,5)^{T}$\\
$(5,2,(8,8))$ & $(5,0)^{R}$, $(8,1)^{R}$, $(8,6)^{C}$, $(6,7)^{C}$, $(4,8)^{T}$, $(1,7)^{T}$, $(0,5)^{T}$\\
$(7,3,(11,12))$ & $(7,0)^{R}$, $(10,2)^{R}$, $(11,5)^{R}$, $(11,12)^{C}$, $(4,12)^{T}$, $(1,10)^{T}$, $(0,7)^{T}$\\
$(7,2,(9,13))$ & $(7,0)^{R}$, $(8,2)^{R}$, $(9,4)^{R}$, $(9,11)^{C}$, $(6,12)^{T}$, $(3,13)^{T}$, $(2,11)^{T}$, $(1,9)^{T}$, $(0,7)^{T}$\\
$(8,3,(12,10))$ & $(8,0)^{R}$, $(11,1)^{R}$, $(12,4)^{R}$, $(11,9)^{C}$, $(6,10)^{T}$, $(3,9)^{T}$, $(0,8)^{T}$\\
\bottomrule\end{tabular}
\caption{The outer chain $\Omega$ (from $B$ counterclockwise to $A$) with classes as superscripts.}\label{tab:omega-app}
\end{table}

Over all nonspecial classes with $5\le d\le8$ and all boxes $Q=(d+u,d+v)$, $0\le u,v<d$, $(u,v)\ne(0,0)$ (a total of $\statTotal$ boxes), the word of classes of $\Omega$ lies in $R^*C^*T^*$ in $\statWord$ boxes, and the point $C_0=(d,d)$ is a strictly interior point of $H_Q$ in $\statInterior$ boxes (computed by comparing the distance from $C_0$ to the boundary of the hull with $10^{-9}$). This is the numerical content of the correction D9: the position of $C_0$ depends on the box, and it is generally not on the outer chain.

%% -------------------------------------------------------------------
\section{Cap coverage: the cell family in an example}\label{app:cap}

For $d=7$, $s=3$ and $Q=(11,12)$ ($u=4$, $v=5$) the explicit cell family of Section~\ref{sec:cellfamily} has the members listed in Table~\ref{tab:family-app}; the picture is Figure~\ref{fig:sweeps}.

\begin{table}[ht]
\centering\small
\begin{tabular}{@{}llp{8.2cm}@{}}
\toprule family & number & corner 4-tuples $(w,w+g,w+g+h,w+h)$\\\midrule
L-cells & 2 & $(10,2),(11,5),(4,5),(3,2)$; $(11,5),(11,12),(4,12),(4,5)$\\
D-cells & 2 & $(11,12),(4,12),(4,5),(11,5)$; $(4,12),(1,10),(1,3),(4,5)$\\
bottom cell & 1 & $(6,4),(7,0),(10,2),(9,6)$\\
left cell & 1 & $(2,6),(0,7),(1,10),(3,9)$\\
hole-filling & 4 & $(0,7),(2,6),(7,7),(5,8)$; $(5,1),(5,8),(7,7),(7,0)$; $(7,0),(6,4),(7,7),(8,3)$; $(1,3),(8,3),(7,7),(0,7)$\\
\bottomrule\end{tabular}
\caption{The cap cell family for $d=7$, $s=3$, $Q=(11,12)$, as corner four-tuples. Together with the cells of $U_0$ these cells cover $H_Q$ (numerical test).}\label{tab:family-app}
\end{table}

Table~\ref{tab:records-app} shows the record minima, the generator $g$ chosen in the bottom cell, and the verification of the inequality $\cross(g,M)\ge xy$ of Lemma~\ref{lem:arith} for $d=13$, $s=5$ and $u=1,\dots,12$.

\begin{table}[ht]
\centering\small
\begin{tabular}{@{}rrrcc@{}}
\toprule $u$ & record $M=(x,y)$ & $g$ & $\cross(g,M)$ & $xy$\\\midrule
1 & $(1,5)$ & $(1,-8)$ & 13 & 5\\
2 & $(1,5)$ & $(1,-8)$ & 13 & 5\\
3 & $(3,2)$ & $(1,-8)$ & 26 & 6\\
4 & $(3,2)$ & $(1,-8)$ & 26 & 6\\
5 & $(3,2)$ & $(1,-8)$ & 26 & 6\\
6 & $(3,2)$ & $(1,-8)$ & 26 & 6\\
7 & $(3,2)$ & $(1,-8)$ & 26 & 6\\
8 & $(8,1)$ & $(1,-8)$ & 65 & 8\\
9 & $(8,1)$ & $(1,-8)$ & 65 & 8\\
10 & $(8,1)$ & $(1,-8)$ & 65 & 8\\
11 & $(8,1)$ & $(1,-8)$ & 65 & 8\\
12 & $(8,1)$ & $(1,-8)$ & 65 & 8\\
\bottomrule\end{tabular}
\caption{Records of $5i\bmod13$ ($i\le u$), the chosen Graver vector $g$ and the inequality of Lemma~\ref{lem:arith}.}\label{tab:records-app}
\end{table}

%% -------------------------------------------------------------------
\section{Computational certification}\label{app:comp}

\subsection{What is in the package}

\begin{description}
\item[\texttt{scripts/geo.py}] exact-integer geometry helpers (fibres, Graver vectors, cells, boundary chain) with \texttt{shapely} used only for areas and unions.
\item[\texttt{scripts/figs.py}] generates the twelve figures and \texttt{figures/figure\_data.json}.
\item[\texttt{scripts/realize.py}] the realizability check of Proposition~\ref{prop:realize} for $2\le d\le14$ (primes are constructed as in the proof; the congruence lattice is compared with $\Gamma_{d,s}$ over a box).
\item[\texttt{scripts/example.py}] the exact worked example of Section~\ref{sec:example} (integers and \texttt{Fraction} only).
\item[\texttt{scripts/verify\_all.py}] the finite tests of Table~\ref{tab:checks}, in two modes: \texttt{smoke} (small ranges, about six seconds) and \texttt{full}.
\item[\texttt{scripts/gen\_tables.py}] regenerates the data tables of the appendices.
\item[\texttt{logs/}] the logs of the runs reported in this paper.
\end{description}

\subsection{Environment}

The runs reported were made with Python 3.13, \texttt{shapely} 2.2.0, \texttt{matplotlib} 3.11.2 and pdfTeX 3.141592653-2.6-1.40.25 (TeX Live 2023); these versions are pinned in \texttt{requirements.txt}. The pinned GEO-I source commit is \geoicommit. The \LaTeX{} build has no errors, no warnings and no overfull boxes; the log contains underfull-box notices (30 in our environment), which are typographic diagnostics and do not affect content. All scripts use only the standard library apart from \texttt{shapely} and \texttt{matplotlib}.

\subsection{Results}

\begin{table}[ht]
\centering\small
\begin{tabular}{@{}lrrrr@{}}
\toprule
Test & smoke instances & failures & full instances & failures\\
\midrule
cap cover ($5\le d\le8$ / $\le12$) & 366 & 0 & 2130 & 0\\
Pareto inequality ($d\le40$ / $\le120$) & 2114 & 0 & 33\,633 & 0\\
boundary steps ($d\le8$ / $\le12$) & 997 & 0 & 2343 & 0\\
universal box test ($d\le6$ / $\le8$) & 451 & 0 & 1327 & 0\\
\bottomrule
\end{tabular}
\caption{Results of \texttt{verify\_all.py} in the two modes.}\label{tab:verify-app}
\end{table}

\subsection{Nature of the checks}

The geometric checks form unions of polygons with \texttt{shapely} in double precision and compare areas with a tolerance of $10^{-9}$. They can in principle report a missing sliver of area below the tolerance as covered. Because all vertices have small integer coordinates and the candidate gaps are triangles of area at least $\tfrac12$, such a case would not be missed by more than a numerical margin; but this is an argument about the instances, and the checks remain numerical tests. The arithmetic checks (Pareto inequality, realizability, boundary steps, worked example) are in exact integer or rational arithmetic.

\subsection{Reproducing}

In the package directory, \texttt{python3 scripts/verify\_all.py smoke} reproduces the first two columns of Table~\ref{tab:verify-app} and \texttt{python3 scripts/verify\_all.py full} the last two; \texttt{python3 scripts/example.py}, \texttt{python3 scripts/realize.py} and \texttt{python3 scripts/gen\_tables.py} reproduce the other items; the checksum of every file of the package is in \texttt{MANIFEST.sha256}, checkable with \texttt{sha256sum -c}.

%% -------------------------------------------------------------------
\section{Theorem dependency registry}\label{app:registry}

Table~\ref{tab:registry} relates every result of the paper to its identifier in the GEO-I series. The GEO-I identifiers refer to the archived releases at the pinned commit \geoicommit, with release, document and role given in Tables~\ref{tab:geoi-roles} and~\ref{tab:geoi-theorems}. ``Source status'' is the status recorded in the release; ``here'' is the status of the proof written in this paper (all proofs are written in full; none has been refereed or formalized, Section~\ref{sec:tiers}). The crosswalk used in preparing this table (entries X01--X24) is distributed with the package as \texttt{PAPER\_A\_CROSSWALK.tsv}.

\begin{longtable}{@{}p{2.9cm}p{2.8cm}p{3.0cm}p{4.6cm}@{}}
\toprule Manuscript & GEO-I ID & Source status & Status here\\\midrule\endhead
Prop.~\ref{prop:bij}, Thm~\ref{thm:occ} & v001 Prop.~1.2, Thm~1.4 & proved; checked samples & proved in Section~\ref{sec:arith}\\
Prop.~\ref{prop:norm}, \ref{prop:transfer} & v004 Prop.~S1 & proved & proved in Section~\ref{sec:planar}; $k=2$ only\\
Prop.~\ref{prop:realize} & --- & not in releases & new in this manuscript; proof + finite test; nothing depends on it\\
Prop.~\ref{prop:elem}, Def.~\ref{def:cells} & v004/v005/v012 \S0 & definitions; facts asserted & definitions restated; facts proved\\
Lemma~\ref{lem:special-graver} & v005 Lemma~G & proved & proved in Section~\ref{sec:special-stab}\\
Thm~\ref{thm:L} & v005 Thm~L = G5-C1 & proved + checked & proved; range $Q\ge(d,d)$, $Q\ne(d,d)$ (D1)\\
Prop.~\ref{prop:special-fan} & v005 G5-C2 ($s=-1$) & proved + checked & proved\\
Prop.~\ref{prop:special-plus} & v005 G5-C2 ($s=+1$, $\beta_1$ only) & $\beta_1=2$ proved + checked & exact complement proved here (D3), not attributed to v005\\
Prop.~\ref{prop:hole-exclusion} & G7-T3 (v007 H1--H3) & proved + checked & proved\\
Lemmas~\ref{lem:chain0}--\ref{lem:cover0}, Thm~\ref{thm:exceptional} & G10-L1, G10-Z, G10-T1 (= G9-C1) & proved + checked & proved\\
Thm~\ref{thm:betti} (nonspecial) & G10-T1H (= G9-T1H) & proved + checked & proved over $\mathbb Q$; Alexander duality cited\\
Prop.~\ref{prop:fill} & G8-T1 & proved + checked & proved with explicit cells $R_1,R_2$\\
Cor.~\ref{cor:old} & G8-R1 (conditional), G10-R1 & proved (conditional) / proved + checked & unconditional via Thm~\ref{thm:exceptional} (D5)\\
Thm~\ref{thm:large} & G7-T1 (A-COL) & proved + checked & proved in range $Q\ge(d,d)$, $\max Q\ge2d$\\
Lemma~\ref{lem:arith} & G12-T1 (= G11-A) & proved + checked & proved, nonspecial scope explicit\\
Lemma~\ref{lem:reflection} & G12-T2 & proved + checked & proved for all $Q\ge(d,d)$ (D4)\\
Prop.~\ref{prop:persist}, Lemmas~\ref{lem:omega}, \ref{lem:records}, \ref{lem:last} & G12-L1, L4, L5 (L2, L3 unregistered) & proved + checked & proved (D6, D9)\\
Lemmas~\ref{lem:L}--\ref{lem:bottom}, Thm~\ref{thm:cap} & G12-T3..T7 & proved + checked (T6 negative) & proved; no residual pentagon\\
Thm~\ref{thm:universal} & G12-T8 & proved + checked, from G5-C1, G7-T1, G10-R1 & assembled from proved results in this paper\\
Cor.~\ref{cor:actual} & v004 S1 + G12-T8 & --- & derived here\\
Question~\ref{q:global} & G12-O1 & open & open\\
\bottomrule
\caption{Registry of results.}\label{tab:registry}
\end{longtable}

\subsection{Dependency structure}

The logical dependencies among the main results are as follows.
\begin{itemize}
\item Theorem~\ref{thm:universal} $\Leftarrow$ Theorem~\ref{thm:L} (special classes); Theorem~\ref{thm:large} (a side $\ge2d$); Corollary~\ref{cor:old} and Theorem~\ref{thm:cap} (the remaining small boxes).
\item Theorem~\ref{thm:L} $\Leftarrow$ Lemmas~\ref{lem:special-graver} and~\ref{lem:slices}.
\item Theorem~\ref{thm:large} $\Leftarrow$ Lemma~\ref{lem:columns} and Lemma~\ref{lem:reflection}.
\item Corollary~\ref{cor:old} $\Leftarrow$ Theorem~\ref{thm:exceptional}, Proposition~\ref{prop:fill}, Proposition~\ref{prop:elem}.
\item Theorem~\ref{thm:exceptional} $\Leftarrow$ Proposition~\ref{prop:hole-exclusion}, Lemmas~\ref{lem:chain0}, \ref{lem:zonogon}, \ref{lem:K}, \ref{lem:cover0}.
\item Theorem~\ref{thm:cap} $\Leftarrow$ Proposition~\ref{prop:persist}, Lemmas~\ref{lem:omega}, \ref{lem:records}, \ref{lem:last}, \ref{lem:L}, \ref{lem:D}, \ref{lem:bottom}, Lemma~\ref{lem:arith}, Lemma~\ref{lem:reflection}.
\item Lemma~\ref{lem:reflection} $\Leftarrow$ Lemma~\ref{lem:axis-pair}.
\item Theorem~\ref{thm:betti} $\Leftarrow$ Theorem~\ref{thm:exceptional}, Propositions~\ref{prop:special-fan}, \ref{prop:special-plus}, Corollary~\ref{cor:planar-betti}.
\item Corollary~\ref{cor:actual} $\Leftarrow$ Propositions~\ref{prop:norm}, \ref{prop:transfer}, Corollary~\ref{cor:plateau}, Theorems~\ref{thm:universal}, \ref{thm:betti}.
\end{itemize}
Proposition~\ref{prop:realize} is not an ancestor of any result.
